\documentclass[11pt,a4paper]{article}

\usepackage{../../shared/cathedral-preamble}
\usepackage{setspace}
\onehalfspacing

\title{%
  The Dark Mirror:\\
  Dual-Use Risks of Number-Theoretic Engineering\\
  from the Cathedral Project%
}
\author{
  Jason Robert Gochanour
}
\date{June 2026}

\begin{document}

\maketitle

\begin{abstract}
The companion papers \emph{What Could Be Built} and \emph{The
Prime Number Grid} propose engineering applications of the
mathematical structures discovered during the Cathedral project.
Responsible disclosure requires examining the same technologies
from an adversarial perspective. This paper identifies seven
dual-use risks: (1)~precision harmonic injection attacks on power
grids, (2)~resonance-targeted industrial sabotage via arithmetic
vibration analysis, (3)~adversarial use of verified instability
certificates for cascade failure engineering, (4)~broadband stealth
and jamming from prime-spaced antenna arrays, (5)~covert RF
surveillance via arithmetic energy harvesting;
(6)~AI-accelerated mathematical attack planning, where large
language models lower the expertise barrier from graduate-level
mathematics to natural-language prompting; and
(7)~algorithmic blinding via arithmetic denial-of-service,
where high-condition-number spectral injections cause
defender software to hallucinate.
Since the original version of this paper (April~20, 2026),
the Cathedral has evolved from a single-path, single-axiom
architecture (v5) to a \emph{two-axiom crown} architecture (v23) with
multiple independent compiler-verified proof routes connected by the
Parseval Bridge---a zero-axiom isometry between time and frequency
domains---and a Torus Projection that decomposes the Gram energy onto
the infinite-dimensional torus $T^\infty = \prod_p S^1_p$.
Four universal engineering design patterns (Cholesky cooling,
Ward Health Index, anomaly-gated precision, torus-stratified
parallelism) create new dual-use surface area that this
updated edition addresses.
This architectural upgrade strengthens both the defensive
and offensive implications of formal verification.
For each risk, we
assess the severity, the novelty (whether the attack vector is
genuinely new), the current state of defenses, and recommended
mitigations. We conclude that the mathematical structures in the
Cathedral do not create fundamentally new attack capabilities, but
could make existing attack vectors more \emph{precise} and harder
to detect. The most significant new development is the
demonstration that AI pair programming can produce formally verified
mathematical infrastructure---lowering the barrier to both
defensive and offensive applications.
This paper exists because the authors believe
that \emph{thinking about misuse before deployment} is a moral
obligation of anyone who builds powerful tools.
\end{abstract}

\tableofcontents

% ================================================================
\section{Why This Paper Exists}
% ================================================================

Every powerful tool has a dual use. The same chemistry that produces
fertilizer produces explosives. The same physics that enables
nuclear medicine enables nuclear weapons. The same number theory
that could improve power grid monitoring could, in principle, be
turned against the grid.

The companion papers propose six positive energy applications. This
paper examines what happens when each is reflected through an
adversarial lens. We do not provide implementation details for any
attack---we describe the \emph{category} of risk and the
\emph{structural} vulnerability, not the recipe.

\subsection{Responsible Disclosure Rationale}

We publish this analysis openly for three reasons:

\begin{enumerate}
  \item \textbf{Attackers already know the math.} The M\"obius
    function, Fourier analysis, spectral theory, and variational
    methods are textbook material. Suppressing this paper would not
    deny attackers any capability. Moreover, the emergence of AI
    mathematical assistants means that the effective barrier to
    applying these tools is no longer a graduate education---it is
    the ability to describe the problem to a language model.

  \item \textbf{Defenders need to understand the threats.} Grid
    operators, infrastructure security teams, and turbine
    manufacturers benefit from understanding how number-theoretic
    tools could be used against their systems. Silence helps
    attackers more than defenders.

  \item \textbf{Moral obligation.} If we propose technologies, we
    must also assess their misuse potential. Publishing only the
    positive applications would be intellectually dishonest and
    ethically incomplete.
\end{enumerate}

\subsection{What Changed: The Dual Path Architecture}

Since the initial version of this paper (April~20), the Cathedral
has undergone a fundamental architectural transformation:

\begin{itemize}
  \item \textbf{One-Pillar Architecture}: The crown theorem now depends
    on exactly one literature axiom (\texttt{baez\_duarte\_forward},
    IMRN 2003). Three alternative proof paths (Mellin, Perron,
    Renormalization) are preserved with different axiom footprints,
    all connected by the Parseval Bridge (zero-axiom isometry).
  \item \textbf{Scale}: 390+ active Lean files, $\sim$120{,}000 lines of
    formally verified logic, 55+ computational experiments,
    0 sorry placeholders on the critical path, 8{,}736 build jobs.
  \item \textbf{DD-precision certification}: A Dekker--Knuth
    Double-Double Conjugate Gradient engine certifies the
    B\'aez-Duarte distance at $N = 55{,}440$ ($\dim = 55{,}439$),
    the largest such computation in the literature. A 512-bit MPFR
    engine provides Sherman-Morrison cross-validation to $10^{-17}$.
  \item \textbf{AI collaboration}: The entire v5 $\to$ v23 journey
    was conducted through human-AI pair programming over 68 days,
    demonstrating the feasibility of AI-assisted formal verification.
  \item \textbf{Engineering design patterns (v23)}: Four universal
    patterns now extend the attack surface beyond spectral analysis:
    Cholesky cooling (iterative resource allocation), Ward Health
    Index (parity diagnostics), anomaly-gated precision (arithmetic
    vs.\ archimedean separation), and torus-stratified parallel
    computation ($T^\infty$).
\end{itemize}

These changes affect the dual-use landscape. Each risk section
below notes where the v17 architecture alters the assessment.

% ================================================================
\section{Risk 1: Precision Harmonic Injection Attacks on Power Grids}
\label{sec:harmonics}
% ================================================================

\subsection{The Positive Technology}

The M\"obius Harmonic Analyzer (MHA) decomposes grid harmonics into
independent sources, enabling targeted active filtering.

\subsection{The Dark Mirror}

The same M\"obius decomposition that identifies harmonic sources
can be used to \emph{synthesize} optimal harmonic injections. An
attacker with a sufficiently powerful inverter connected to the
grid could:

\begin{enumerate}
  \item \textbf{Targeted transformer saturation}: Inject
    even-order harmonics (2nd, 4th) that cause DC offset in the
    transformer core flux, leading to half-cycle saturation.
    The M\"obius decomposition identifies exactly which harmonic
    frequencies will bypass existing passive filters.

  \item \textbf{Capacitor bank resonance amplification}: Identify
    the resonant frequency of grid capacitor banks
    ($f_r = 1/(2\pi\sqrt{LC})$) and inject power at that frequency.
    If $f_r$ is near an integer harmonic, the arithmetic structure
    of the injection maximizes energy transfer.

  \item \textbf{Protective relay confusion}: Craft a waveform whose
    zero-crossing pattern mimics a fault condition, causing
    protective relays to trip unnecessarily. The M\"obius-optimal
    harmonic mixture minimizes the injected power needed to trigger
    a trip.

  \item \textbf{Stealth}: By using M\"obius inversion, the attacker
    can design an injection that appears as ``natural'' harmonic
    distortion (consistent with normal nonlinear load behavior)
    rather than an obvious attack signal.
\end{enumerate}

\subsection{Severity Assessment}

\begin{center}
\begin{tabular}{@{}ll@{}}
\toprule
\textbf{Factor} & \textbf{Assessment} \\
\midrule
Novelty & Low --- harmonic injection is a known attack vector \\
Precision gain & \textbf{High} --- Parseval Bridge (v15) provides \\
Stealth gain & \textbf{High} --- attack looks like natural distortion \\
Physical access required & Grid connection (industrial inverter) \\
Detection difficulty & High if designed using M\"obius structure \\
Potential impact & Localized transformer/capacitor damage \\
\bottomrule
\end{tabular}
\end{center}

\subsection{Defenses}

\begin{enumerate}
  \item \textbf{Deploy the MHA defensively}: The same tool that
    enables the attack also detects it. An MHA-equipped grid monitor
    can identify ``unnatural'' harmonic source patterns that don't
    match known load profiles.
  \item \textbf{Harmonic injection limits}: IEEE~519 already sets
    limits on harmonic current injection. Enforce these at the
    point of common coupling with real-time monitoring.
  \item \textbf{Anomaly detection}: Train ML models on the
    M\"obius-decomposed harmonic spectrum to detect injections that
    deviate from historical patterns.
\end{enumerate}

\textbf{v15 update}: The Parseval Bridge
(\texttt{White/Scattering.lean}, proved, 0 axioms) now provides a
\emph{machine-verified} bijection between time-domain waveforms and
frequency-domain spectral representations. An attacker could use this
to compute the $L^2$-optimal injection waveform for a given target
spectrum, with verified energy conservation. However, the same bridge
gives defenders an equally powerful detection tool: any anomalous
frequency pattern can be immediately converted to its unique
time-domain signature, making stealth harder, not easier.

% ================================================================
\section{Risk 2: Resonance-Targeted Industrial Sabotage}
\label{sec:vibration}
% ================================================================

\subsection{The Positive Technology}

The Arithmetic Vibration Monitor (AVM) uses compressed sensing with
an arithmetic dictionary to detect incipient turbine faults.

\subsection{The Dark Mirror}

If arithmetic compressed sensing can detect subtle vibration
signatures, it can also \emph{design} them. An attacker with
physical access to a turbine (or the ability to modulate its load)
could:

\begin{enumerate}
  \item \textbf{Resonance excitation}: Identify the natural
    frequencies of turbine blades via the Gram matrix eigenvalue
    structure, then apply periodic load variations at a blade
    natural frequency to induce resonant fatigue failure.

  \item \textbf{Sub-harmonic injection}: Excite blade vibration at
    a frequency that is $1/k$ of a natural frequency ($k$ prime),
    so the vibration builds gradually through nonlinear coupling
    rather than direct resonance. This is harder to detect because
    the excitation frequency is not near any monitored harmonic.

  \item \textbf{Monitoring evasion}: Design the excitation signal
    to be orthogonal to the AVM's sensing dictionary, making it
    invisible to the monitoring system while still delivering
    destructive energy to the blade.

  \item \textbf{Ward Identity Spoofing (v23):} The Ward
    Identity ($v^T G v = D + W$, where $D$ is diagonal self-energy
    and $W$ is off-diagonal Ward current) enables a lightweight
    ``Ward Health Index'' checksum for real-time integrity
    monitoring. An attacker can calculate the exact diagonal
    self-energy required to offset a malicious off-diagonal
    vibration. By injecting a perfectly balanced $(D, W)$
    perturbation that satisfies the Ward Identity, the attacker
    turns off the defender's checksum while the system shakes
    itself to pieces.
\end{enumerate}

\subsection{Severity Assessment}

\begin{center}
\begin{tabular}{@{}ll@{}}
\toprule
\textbf{Factor} & \textbf{Assessment} \\
\midrule
Novelty & Medium --- resonance attacks exist, arithmetic targeting is new \\
Precision gain & \textbf{High} --- spectral analysis identifies exact modes \\
Stealth gain & \textbf{High} --- sub-harmonic excitation evades monitoring \\
Physical access required & Load modulation or physical access to nacelle \\
Detection difficulty & Very high for sub-harmonic approach \\
Potential impact & Blade failure, turbine destruction, potential casualties \\
\bottomrule
\end{tabular}
\end{center}

\subsection{Defenses}

\begin{enumerate}
  \item \textbf{Broadband monitoring}: Monitor vibration across a
    wider frequency range than just the expected harmonics.
    Specifically monitor sub-harmonics at $f_{\text{nat}}/p$ for
    prime $p$.
  \item \textbf{Load pattern anomaly detection}: Flag unusual load
    variation patterns, especially periodic patterns at non-integer
    frequencies.
  \item \textbf{Physical security}: Restrict physical access to
    turbine nacelles and control systems. This is the most effective
    defense against all vibration-based attacks.
\end{enumerate}

% ================================================================
\section{Risk 3: Verified Instability Certificates}
\label{sec:instability}
% ================================================================

\subsection{The Positive Technology}

The Cathedral's framework produces machine-verified stability
certificates for grid configurations: formal proofs that a
Lyapunov function is positive definite.

\subsection{The Dark Mirror}

The same framework can produce machine-verified \emph{instability}
certificates: formal proofs that a specific perturbation sequence
causes cascade failure.

\begin{enumerate}
  \item \textbf{Optimal cascade sequences}: Given a grid model,
    find the minimum set of generator trips that causes voltage
    collapse. The Lean~4 verification ensures the attack plan is
    mathematically correct --- no ``wasted'' disruptions.

  \item \textbf{Stability boundary mapping}: Identify the exact
    operating conditions where the grid transitions from stable to
    unstable. An attacker who can push the grid to this boundary
    (e.g., by coordinating large loads) can trigger instability with
    minimal visible action.

  \item \textbf{Targeted N-k analysis}: Current grid security
    standards require N-1 or N-2 contingency analysis (survive the
    loss of 1 or 2 components). A verified instability certificate
    could identify the specific N-3 or N-4 scenarios that current
    analysis misses.
  \item \textbf{Composite Anchor Targeting (v15):} Recent spectral
    telemetry reveals exactly where systemic instability localizes.
    The ground state of the spatial Hamiltonian---the mathematical
    vector of collapse---exhibits quantum scarring: it avoids
    prime-indexed (low-connectivity) nodes, instead localizing
    onto highly composite numbers (massive, highly connected
    topological heatsinks). An attacker armed with a verified
    instability certificate can target the specific composite
    anchors that hold a system's spectral gap open, rather than
    searching the entire topology.

  \item \textbf{Equation of State Chaos Triggering (v15):}
    Cathedral experiments demonstrate that prime-indexed systems
    transition into uncontrollable GOE chaos at a precisely
    predictable threshold: $N_c \approx 60 \times (m/\varphi(m))$.
    An adversary can quietly perturb a system to $N_c - 1$, where
    it appears mathematically stable, then trigger a single
    additional perturbation that snaps the system into chaotic
    thermalization, overwhelming all linear response controllers.

  \item \textbf{Torus-Stratified Targeting (v23):}
    The Torus Projection ($T^\infty = \prod_p S^1_p$) decomposes
    system energy onto per-prime circles. An attacker can identify
    which prime strata carry the most system energy (e.g., $p = 2$
    carries 46--86\% in the Cathedral's measurements) and
    target \emph{that specific stratum}---disabling the dominant
    harmonic while leaving the system superficially operational.
    This is analogous to targeting a specific quark color in QCD
    rather than attacking the whole proton.

  \item \textbf{Cholesky Decrement Exploitation (v23):}
    The identity $d^2_{N+1} = d^2_N - y^2_{\text{new}}$ reveals
    the \emph{exact marginal contribution} of each system component.
    An attacker can compute the Schur complement
    $y^2_k = (b_k - g^T G_{-k}^{-1} b)^2 / (G_{kk} - g^T G_{-k}^{-1} g)$
    for every node $k$, producing a rank-ordered list of the most
    critical components. Removing the node with the largest $y^2_k$
    causes the maximum increase in system residual---mathematically
    optimal sabotage.
\end{enumerate}

\subsection{Severity Assessment}

\begin{center}
\begin{tabular}{@{}ll@{}}
\toprule
\textbf{Factor} & \textbf{Assessment} \\
\midrule
Novelty & \textbf{High} --- verified attack plans are new \\
Precision gain & \textbf{Very High} --- mathematically optimal attack \\
Stealth gain & Medium --- the attack itself is visible \\
Access required & Grid topology and parameters (often semi-public) \\
Detection difficulty & Low (the cascade is observable) \\
Potential impact & \textbf{Regional blackout, economic damage} \\
\bottomrule
\end{tabular}
\end{center}

This is the most concerning risk in this paper. The novelty is not
in the attack (deliberate grid disruption is well-studied) but in
the \textbf{verification}: a machine-checked proof that the attack
\emph{will} succeed, removing uncertainty from the attacker's
planning.

\textbf{v15 update}: The dual-path architecture significantly amplifies
this risk. A verified instability certificate can now be confirmed
through \emph{two independent mathematical frameworks}: the Mellin
Crown (continuous frequency-domain analysis) and the Gram Crown
(discrete, covariance-free algebraic bounding). If both paths agree that a
perturbation sequence causes instability, the confidence is
qualitatively different from a single-path verification---it is
analogous to the difference between a single witness and two
independent witnesses in law. The Parseval Bridge proves the
mathematical equivalence of both frameworks, meaning an instability
certificate verified in one domain is \emph{guaranteed} to hold in
the other.

Moreover, the B\'aez-Duarte numerical certification engine (512-bit
MPFR, Sherman-Morrison cross-validation to $10^{-17}$) provides a
third, independent validation layer. An attacker who validates their
instability model against the BD engine's numerical output gains
quantitative confidence that their formal proof corresponds to
physical reality.

\subsection{Defenses}

\begin{enumerate}
  \item \textbf{Use the same tool defensively}: If the attacker can
    find verified instabilities, the grid operator should find them
    first. Deploy the same Lean~4 framework to systematically
    identify and remediate all N-3 and N-4 vulnerabilities.
    \textbf{Critically, require dual-path verification}: any
    stability certificate should be verified through both the Mellin
    and Spatial paths to achieve cryptographic-grade trust.
  \item \textbf{Protect grid models}: Limit public access to
    detailed grid topology and impedance data. Current CEII
    (Critical Energy Infrastructure Information) rules should be
    reviewed and strengthened.
  \item \textbf{Dynamic reconfiguration}: Design grids that can
    rapidly change topology in response to detected instability,
    invalidating pre-computed attack plans.
  \item \textbf{Adversarial stability margins}: Design with larger
    stability margins specifically in regions where N-3/N-4
    analysis reveals verified instabilities.
  \item \textbf{Continuous numerical monitoring}: Deploy BD-style
    convergence monitors that continuously validate stability
    margins against the theoretical $C/\ln(N)$ decay rate. Any
    deviation from the predicted trajectory signals either model
    error or active manipulation.
\end{enumerate}

% ================================================================
\section{Risk 4: Prime-Spaced Stealth and Jamming}
\label{sec:stealth}
% ================================================================

\subsection{The Positive Technology}

Prime-spaced sensor arrays with Selberg beamforming weights offer
broadband performance with predictable sidelobe control.

\subsection{The Dark Mirror}

The same array design, used in reverse, can produce electromagnetic
emissions that are:

\begin{enumerate}
  \item \textbf{Broadband noise with minimal power}: A prime-spaced
    jammer distributes its energy across frequencies without
    periodic structure, making it harder to filter. Conventional
    interference cancellation relies on identifying the jammer's
    spectral structure; a prime-spaced jammer has no periodic
    structure to identify.

  \item \textbf{Difficult to geolocate}: The aperiodic radiation
    pattern of a prime-spaced array produces a complex spatial
    pattern that conventional direction-finding algorithms
    (designed for uniform or periodic arrays) may fail to resolve.

  \item \textbf{Radar-evasive scattering}: A target coated in
    prime-spaced resonant elements would scatter radar returns
    across frequencies without the signature Bragg peaks of a
    periodic structure, reducing radar cross-section.
\end{enumerate}

\subsection{Severity Assessment}

\begin{center}
\begin{tabular}{@{}ll@{}}
\toprule
\textbf{Factor} & \textbf{Assessment} \\
\midrule
Novelty & Low --- aperiodic jamming/stealth are well-studied \\
Precision gain & Low --- random spacing works comparably \\
Stealth gain & Low --- prime spacing has marginal advantage \\
Access required & RF hardware (commercially available) \\
Detection difficulty & Medium \\
Potential impact & Communications disruption \\
\bottomrule
\end{tabular}
\end{center}

\textbf{Honest assessment}: This is the weakest risk in the paper.
Prime-spaced arrays are marginally better than random arrays for
jamming/stealth purposes, and the advantage does not justify
concern beyond existing aperiodic jamming threats.

\textbf{v15 update}: However, the discovery of the arithmetic-GOE
constant $\alpha \approx 0.47$ introduces a subtlety. An attacker can
synthesize jamming signals that perfectly mimic the $\alpha \approx 0.47$
thermodynamic background of natural physical systems, potentially
bypassing AI anomaly detectors trained on purely random noise.

\subsection{Defenses}

Existing anti-jamming techniques (spread spectrum, frequency
hopping, adaptive beamforming) are effective against aperiodic
jammers regardless of whether the aperiodicity comes from random
or prime spacing.

\textbf{v15 addition --- The Arithmetic-GOE Fingerprint:}
Recent physical telemetry reveals that prime-spaced systems do not
generate perfectly random noise; they maintain persistent partial
localization at a participation ratio of exactly $\alpha \approx 0.47$.
Defenders do not need to identify specific frequencies.
By running real-time participation ratio analysis on ambient RF noise,
defenders can detect this unforgeable mathematical signature: if the
background noise suddenly shifts to $\alpha \approx 0.47$, the stealth
system's arithmetic structure is mathematically fingerprinted.

% ================================================================
\section{Risk 5: Covert Broadband RF Surveillance}
\label{sec:surveillance}
% ================================================================

\subsection{The Positive Technology}

Prime-indexed rectennas harvest ambient RF energy across a wide
bandwidth, enabling battery-free IoT sensors.

\subsection{The Dark Mirror}

A device that harvests broadband RF also \emph{receives} broadband
RF. A prime-indexed rectenna array disguised as an energy harvester
could:

\begin{enumerate}
  \item \textbf{Passive SIGINT}: Simultaneously monitor WiFi
    (2.4/5 GHz), cellular (700 MHz--3 GHz), Bluetooth (2.4 GHz),
    and IoT (900 MHz) without active emissions, making it
    undetectable by RF sweeps that look for transmitters.

  \item \textbf{Self-powered}: The same harvested energy that
    powers the device also carries the intelligence. No battery
    means no maintenance visits, no supply chain to trace, and
    indefinite operational life.

  \item \textbf{Plausible deniability}: If discovered, the device
    appears to be a legitimate energy harvester (IoT sensor, smart
    building component) rather than a surveillance device.
\end{enumerate}

\subsection{Severity Assessment}

\begin{center}
\begin{tabular}{@{}ll@{}}
\toprule
\textbf{Factor} & \textbf{Assessment} \\
\midrule
Novelty & Medium --- passive RF surveillance exists, self-powered is new \\
Precision gain & Low \\
Stealth gain & \textbf{High} --- disguised as legitimate IoT hardware \\
Access required & Physical placement in target facility \\
Detection difficulty & \textbf{Very high} --- no RF emissions to detect \\
Potential impact & Espionage, corporate/government intelligence \\
\bottomrule
\end{tabular}
\end{center}

\subsection{Defenses}

\begin{enumerate}
  \item \textbf{Physical inspection}: Regular RF environment audits
    that include physical inspection of all connected devices.
  \item \textbf{Faraday shielding}: Sensitive facilities should
    control RF ingress/egress regardless of harvesting technology.
  \item \textbf{Supply chain security}: Verify the provenance of
    all IoT/energy harvesting devices deployed in sensitive
    environments.
\end{enumerate}

% ================================================================
\section{Risk 6: AI-Accelerated Mathematical Attack Planning}
\label{sec:ai}
% ================================================================

\subsection{The Positive Technology}

The Cathedral was built through human-AI pair programming: a human
computer scientist directing AI assistants (Claude, Gemini) to formalize,
debug, and verify Lean~4 proofs. This collaboration model reduced the
development time for the v5~$\to$~v17 evolution from an estimated
months of solo work to approximately two weeks.

\subsection{The Dark Mirror}

The same AI-assisted formalization methodology could be applied to
adversarial engineering:

\begin{enumerate}
  \item \textbf{Barrier reduction}: The original version of this paper
    (April~20) identified ``graduate education in mathematics'' as the
    expertise barrier for applying number-theoretic tools offensively.
    AI mathematical assistants substantially lower this barrier. An
    attacker no longer needs to understand the M\"obius function,
    contour integration, or spectral theory---they need to describe
    \emph{what they want to achieve} in natural language, and the AI
    can formalize the mathematical attack plan.

  \item \textbf{Verification assistance}: The same AI tools that
    help a human verify a stability proof can help an adversary verify
    an instability proof. The AI does not distinguish between
    defensive and offensive goals---it formalizes whatever the human
    asks it to formalize.

  \item \textbf{Rapid iteration}: The Cathedral's development
    demonstrates that AI pair programming enables rapid exploration of
    mathematical strategies. An attacker could iterate through
    multiple attack designs in hours rather than weeks, testing each
    against the formal verification framework until finding one that
    produces a verified certificate.

  \item \textbf{Knowledge synthesis}: AI assistants can draw on the
    full corpus of published mathematics, engineering, and security
    literature simultaneously. An attacker who understands grid
    operations but not number theory---or vice versa---can use AI to
    bridge the knowledge gap.

  \item \textbf{Abstract algebraic jailbreak (v23):} Mathematics
    is the ultimate AI safety bypass. If an attacker prompts an
    LLM with ``Write a script to destabilize a power grid,'' every
    commercial AI will trigger a safety refusal. But if the attacker
    prompts ``Help me write a Lean~4 proof identifying the node $k$
    that maximizes the Schur complement $y^2_k$ in the Cholesky
    decomposition of this adjacency matrix,'' the AI will
    enthusiastically comply. Current AI alignment models cannot
    detect malicious intent when it is compiled into pure abstract
    algebra. By translating physical sabotage into a Gram matrix
    problem, the attacker effectively blinds the AI's safety
    filters---the mathematical content is indistinguishable from
    legitimate research.
\end{enumerate}

\subsection{Severity Assessment}

\begin{center}
\begin{tabular}{@{}ll@{}}
\toprule
\textbf{Factor} & \textbf{Assessment} \\
\midrule
Novelty & \textbf{Very High} --- AI-assisted formal attack planning is new \\
Precision gain & \textbf{High} --- AI enables mathematical optimization \\
Stealth gain & Low --- the AI collaboration itself is not stealthy \\
Access required & Internet access and an AI subscription \\
Detection difficulty & N/A (planning phase, not execution) \\
Potential impact & Force multiplier for all other risks \\
\bottomrule
\end{tabular}
\end{center}

\textbf{Honest assessment}: This is the most important new risk
since the paper's original publication. It is also the most
difficult to mitigate, because the same AI capabilities that
enable offensive planning are essential for defensive analysis.
Any restriction on AI-assisted mathematics would harm defenders
(who operate openly and benefit from AI assistance) more than
attackers (who can use unrestricted models or open-source
alternatives).

\subsection{Defenses}

\begin{enumerate}
  \item \textbf{Normalize defensive AI use}: The most effective
    defense is to ensure that infrastructure operators adopt
    AI-assisted formal verification \emph{before} attackers do.
    The defender's advantage is system access and domain knowledge;
    AI amplifies this advantage if the defender uses it first.
  \item \textbf{AI safety alignment}: AI systems should be designed
    to recognize and flag requests that appear to target critical
    infrastructure, while preserving the ability to assist
    legitimate defensive analysis. This is a difficult but
    necessary balance.
  \item \textbf{Institutional capacity}: Invest in training
    infrastructure security teams to use AI-assisted formal
    verification tools. The barrier should be lowered for
    defenders, not raised for everyone.
\end{enumerate}

% ================================================================
\section{Risk 7: Algorithmic Blinding (Arithmetic Denial-of-Service)}
\label{sec:blinding}
% ================================================================

\subsection{The Positive Technology}

Grid monitoring systems and the proposed M\"obius Harmonic Analyzer
rely on real-time linear algebra and eigensolvers to track system
stability.

\subsection{The Dark Mirror}

The geometric frustration of prime-spaced topologies naturally
produces interaction matrices with extreme condition numbers
($\kappa > 10^6$). During Cathedral development (Exploration~19),
standard 64-bit floating-point eigensolvers catastrophically failed
at $N \geq 750$, hallucinating negative eigenvalues due to mantissa
exhaustion---precisely where the ground truth eigenvalue was
$\lambda_{\min} \approx 10^{-6}$.

\begin{enumerate}
  \item \textbf{Precision-wall injection}: An attacker can synthesize
    a physical power-line injection patterned on high-condition-number
    prime topologies. When the defender's SCADA software attempts to
    compute the grid state, the mathematical complexity forces
    catastrophic algorithmic cancellation, silently blinding the
    software to the physical reality of the grid.

  \item \textbf{Eigenvalue hallucination}: The defender's stability
    monitor may report the system as unstable (negative eigenvalues)
    when it is actually stable, causing unnecessary protective relay
    trips---or worse, report stability when the system is failing.

  \item \textbf{Pre-kinetic blinding}: The precision-wall attack can
    be used as a precursor to a physical attack. By first blinding
    the defender's monitoring, the attacker creates a window of
    unmonitored operation during which a kinetic disruption proceeds
    undetected.

  \item \textbf{The Dekker--Knuth Boundary (v16):} The Cathedral's
    DD-precision pipeline (May~2026) discovered the exact numerical
    boundary: at dimension $N > 40{,}000$ (condition number
    $\kappa \sim 10^{15}$), \emph{every} standard \texttt{f64}
    eigensolver produces non-physical results. The matrix appears
    non-positive-definite in \texttt{f64} even though it is provably
    SPD at full precision. An attacker armed with this knowledge can
    \emph{calibrate} an injection to sit precisely at the \texttt{f64}
    precision boundary of the defender's software, ensuring the
    monitoring system fails silently rather than catastrophically.
    The defender sees ``stable'' when the system is collapsing,
    or ``unstable'' when the system is fine. This is no longer
    theoretical: the Cathedral measured the exact eigenvalue at
    which \texttt{f64} representation fails ($\sim 10^{-15}$) and the
    Dekker--Knuth DD representation that fixes it ($\sim 10^{-31}$).

  \item \textbf{Archimedean shadowing (v23):} The Glass Bridge
    identity $G = R + \frac{1}{4}\mathbf{1}\mathbf{1}^T$ separates
    the Gram matrix into a fast, rational arithmetic part ($R$,
    pure GCD structure) and a computationally expensive transcendental
    anomaly ($\Delta_{\text{arch}}$). Defenders running real-time
    SCADA telemetry will likely truncate $\Delta_{\text{arch}}$ to
    save compute cycles, computing only the fast rational part.
    An attacker can encode a destructive resonant injection entirely
    within the archimedean anomaly sector. Because the defender's
    software truncates the anomaly, the attack lives in the
    ``blind spot'' of classical sieve mathematics---invisible to
    the monitoring system.
\end{enumerate}

\subsection{Severity Assessment}

\begin{center}
\begin{tabular}{@{}ll@{}}
\toprule
\textbf{Factor} & \textbf{Assessment} \\
\midrule
Novelty & \textbf{High} --- arithmetic DoS via condition number is new \\
Precision gain & \textbf{High} --- exactly calibrated to f64 mantissa \\
Stealth gain & \textbf{Very High} --- the attack is invisible to the blinded system \\
Access required & Grid connection (power-line injection) \\
Detection difficulty & \textbf{High} --- the monitoring system itself is compromised \\
Potential impact & \textbf{Blinding of critical monitoring infrastructure} \\
\bottomrule
\end{tabular}
\end{center}

\subsection{Defenses}

\begin{enumerate}
  \item \textbf{Arbitrary-precision monitoring}: Industrial monitoring
    systems analyzing prime-spaced harmonic data must transition from
    standard \texttt{f64} to computationally expensive arbitrary-precision
    arithmetic (MPFR) or exact rational solvers when condition numbers
    exceed $10^4$.
  \item \textbf{Condition number monitoring}: Real-time tracking of
    the condition number of the system Gram matrix. A sudden increase
    in $\kappa$ should trigger an automatic switch to high-precision
    solvers and alert operators.
  \item \textbf{Dual-solver cross-validation}: Run both \texttt{f64}
    and MPFR eigensolvers in parallel; any disagreement between them
    indicates either precision loss or adversarial injection.
  \item \textbf{DD-precision monitoring (v16):} Critical infrastructure
    monitoring systems should deploy Double-Double (Dekker--Knuth)
    arithmetic for all spectral computations where the condition number
    may exceed $10^{10}$. The cost is approximately $2\times$ compared
    to \texttt{f64}---a trivial overhead for safety-critical systems.
    The Cathedral's \texttt{cathedral-utils} library provides a
    production-grade implementation~\cite{dekker1971}.
\end{enumerate}

% ================================================================
\section{Risk Assessment Summary}
% ================================================================

\begin{center}
\small
\begin{tabular}{@{}lcccl@{}}
\toprule
\textbf{Risk} & \textbf{Novelty} & \textbf{Severity} &
\textbf{Feasibility} & \textbf{Primary Defense} \\
\midrule
Grid harmonic injection & Low & Medium & High & MHA defensive monitoring \\
Vibration sabotage & Medium & High & Medium & Broadband monitoring \\
Verified instability & \textbf{High} & \textbf{Very High} &
  Medium & Dual-path defensive verification \\
Stealth/jamming & Low & Low-Med & High & $\alpha$-fingerprint detection \\
RF surveillance & Medium & Medium & Medium & Physical/supply chain \\
\textbf{AI attack planning} & \textbf{Very High} & \textbf{High} &
  \textbf{Very High} & \textbf{Normalize defensive AI use} \\
\textbf{Algorithmic blinding} & \textbf{High} & \textbf{Very High} &
  High & DD-precision / dual-solver validation \\
\midrule
Instability: torus-targeting & High & High & Medium & Defensive torus monitoring \\
Instability: Cholesky exploit. & \textbf{High} & \textbf{Very High} & Medium & Classify Schur rankings \\
\bottomrule
\end{tabular}
\end{center}

\subsection{The Headline Risk}

The single most concerning capability is \textbf{Risk~3: Verified
Instability Certificates}. The ability to produce a
\emph{machine-checked proof} that a specific attack sequence will
cause a grid cascade failure is qualitatively different from
simulation-based attack planning. A simulation can be wrong; a
verified certificate cannot (within its axioms).

The mitigation is symmetric: grid operators should use the same
framework to find these vulnerabilities first. The race is between
defensive and offensive verification, and the defender has the
advantage of knowing their own system.

% ================================================================
\section{The Broader Dual-Use Question}
% ================================================================

None of the risks identified in this paper are created by the
Cathedral. The mathematical tools (Fourier analysis, spectral
theory, M\"obius function, variational methods) have been publicly
available for decades to centuries. What the Cathedral contributes
is:

\begin{enumerate}
  \item \textbf{Organization}: The Cathedral assembles these tools
    into a coherent framework. An attacker would otherwise need to
    independently recognize the connections between number theory
    and engineering.
  \item \textbf{Verification}: The machine-verified proof
    methodology could, in principle, be applied to verify attack
    plans rather than proofs. The dual-path architecture (v15)
    makes verified claims carry extraordinary weight: two
    independent mathematical witnesses, connected by a proved
    isometry, cannot both be wrong.
  \item \textbf{Precision}: The number-theoretic structure (primes,
    M\"obius function, Selberg weights) provides optimal or
    near-optimal parameters for certain engineering tasks,
    including adversarial ones.
  \item \textbf{Accessibility}: The demonstration that AI pair
    programming can produce formally verified mathematics lowers
    the effective barrier from ``graduate-level expertise'' to
    ``ability to direct an AI assistant.'' This is the most
    significant change since the paper's original publication.
  \item \textbf{Lattice cryptography implications}: The DD-precision
    Conjugate Gradient solver, designed to navigate ill-conditioned
    55,000-dimensional integer lattices, shares the computational
    profile of the Closest Vector Problem (CVP) central to
    lattice-based post-quantum cryptography (NTRU, Kyber, Dilithium).
    While the Gram matrix has specific arithmetic structure not
    present in cryptographic lattices, the mixed-precision solver
    techniques developed for the Cathedral could potentially
    transfer to CVP solvers. This warrants monitoring by the
    cryptographic community.
\end{enumerate}

The dual-use question is not whether to publish (the math is
already public) but whether to \emph{emphasize the connections}
between number theory and infrastructure engineering. We believe
the answer is yes, because:

\begin{quote}
\emph{The asymmetry in dual-use technology favors the defender.}
\end{quote}

Defenders (grid operators, turbine manufacturers, infrastructure
security teams) have domain knowledge, system access, and
institutional resources. Attackers must operate covertly, with
limited access and imperfect models. AI amplifies the defender's
advantage more than the attacker's, because the defender can
leverage AI with full system access while the attacker cannot.
Providing defenders with a
new mathematical framework strengthens them more than it
strengthens attackers.

% ================================================================
\section{Recommendations}
% ================================================================

\begin{enumerate}
  \item \textbf{Deploy the MHA defensively before it is used
    offensively.} The M\"obius Harmonic Analyzer is the most
    immediately actionable tool in both the positive and negative
    applications. Grid operators should adopt it for monitoring
    before attackers adopt it for injection design.

  \item \textbf{Extend N-k contingency analysis.} Use the Cathedral's
    verified stability framework to systematically identify N-3 and
    N-4 grid vulnerabilities. Remediate them before an adversary
    discovers them independently. \textbf{Require dual-path
    verification} (Mellin + Spatial) for all critical stability
    certificates.

  \item \textbf{Add sub-harmonic channels to vibration monitors.}
    Current turbine monitoring focuses on integer harmonics of the
    rotation frequency. Adding monitoring at $f/p$ for small primes
    $p$ would detect the sub-harmonic excitation attack described
    in Risk~2.

  \item \textbf{Classify verified instability certificates.}
    Machine-checked proofs of grid instability should be treated
    as sensitive information (CEII) and protected accordingly.

  \item \textbf{Invest in AI-assisted defensive verification.}
    Train infrastructure security teams to use AI pair programming
    for formal verification of grid stability. The defender who
    uses AI first gains an asymmetric advantage.

  \item \textbf{Develop AI safety guidelines for infrastructure
    mathematics.} Establish norms for how AI assistants should
    handle requests involving critical infrastructure analysis.
    The goal is not to restrict mathematical capability but to
    ensure responsible use.

  \item \textbf{Publish this paper.} Suppressing dual-use analysis
    protects no one. The mathematical tools are public. The
    infrastructure is public. The only thing that can be private is
    the defender's understanding of their own vulnerabilities---and
    that understanding is strengthened, not weakened, by open
    analysis.

  \item \textbf{Notify the Post-Quantum Cryptography (PQC)
    community.} The extreme-precision, high-dimension linear
    algebra solvers developed for the Cathedral's lattice
    (specifically the DD-CG pipeline in \texttt{cathedral-utils})
    should be shared with cryptographic researchers. The PQC
    community must evaluate whether these techniques alter the
    practical complexity of continuous relaxation steps in
    lattice reduction algorithms (e.g., BKZ) used to attack
    current NIST post-quantum standards (ML-KEM/Kyber,
    ML-DSA/Dilithium).
\end{enumerate}

% ================================================================
\section{Conclusion}
% ================================================================

The Cathedral's mathematics is neutral. The Parseval Bridge connects
time and frequency domains. The M\"obius function decomposes
multiplicative structure. The Rayleigh quotient bounds eigenvalues.
These tools do not care whether they are used to protect a power
grid or to attack one.

What matters is who learns to use them first---and in the age of
AI-assisted mathematics, ``learning'' happens faster than ever before.

The dual-path architecture adds a new dimension to this neutrality.
The proof has a gauge symmetry: the same truth can be witnessed in
multiple independent mathematical frameworks, connected by the Parseval
isometry. This makes verified claims---whether stability or
instability certificates---carry extraordinary epistemic weight.
A claim verified in both gauges cannot be forged, but it also
cannot be dismissed.

The v23 engineering design patterns (Cholesky cooling, Ward Health
Index, anomaly gating, torus stratification) extend this neutrality
into four new domains. The Cholesky decrement that helps a sensor
network designer find optimal placement is the same formula that
helps an attacker find the most critical node to remove. The torus
projection that parallelizes computation also reveals per-prime
vulnerabilities. The mathematics does not choose sides.

This paper exists to ensure that the people who protect critical
infrastructure understand the same mathematics---and the same
AI-assisted methodology---that could threaten it. We believe that
open, honest analysis of dual-use risks---even uncomfortable
analysis---is a more effective defense than silence.

The Cathedral was built to illuminate. This paper points the light
at the shadows.

\begin{quote}
\emph{A tool that can diagnose can also deceive.\\
A proof that certifies stability can also certify its absence.\\
An AI that helps a defender can help an attacker.\\
The mathematics does not choose sides.\\
We must.}
\end{quote}

\vspace{0.5cm}

\begin{thebibliography}{99}

\bibitem{ieee519}
IEEE, \emph{IEEE Std 519-2022: Recommended Practice for Harmonic
Control}, 2022.

\bibitem{nerc}
NERC, \emph{Grid Security Guidelines: Critical Infrastructure
Protection Standards CIP-002 through CIP-014}, 2023.

\bibitem{stuxnet}
R.~Langner, \emph{Stuxnet: Dissecting a Cyberwarfare Weapon},
IEEE Security \& Privacy, 2011.

\bibitem{ceii}
FERC, \emph{Critical Energy Infrastructure Information (CEII)
Regulations}, 18 CFR \S 388.113.

\bibitem{cathedral-v17}
J.R.~Gochanour, \emph{The Cathedral: A Formal Reduction of the
Riemann Hypothesis to Two Axioms via the Nyman--Beurling
Criterion (v23)}, 2026.

\bibitem{cathedral-physics}
J.R.~Gochanour, \emph{The Physics of the Primes: A Dictionary
Between the Cathedral Proof and Quantum Mechanics}, 2026.

\bibitem{dekker1971}
T.J.~Dekker, \emph{A floating-point technique for extending the
available precision}, Numer. Math. \textbf{18} (1971), 224--242.

\end{thebibliography}

\end{document}
